\section{Pointwise support checks with an explicit unlabeled-data cost}
\label{sec:unlabeled-stability-pool}

There is another way to avoid a global support restriction or a
low-complexity family of exceptional regions: use enough unlabeled
inputs to pay for unrestricted variation of the instability sets.
The price is an unlabeled sample size that depends on width and prior
scale. The labeled sample bound remains width independent.

Fix one Gaussian learner, with width $N$, $0<t\le1$, and raw prior
$\mu=\mu_{N,t}$. Put $a=N/t^2$. Let $\mathsf{Stable}(w,x)$ denote
the primary clipped-score condition at the single input $x$, with
its calibrated radius $r_c(w)$. The margin-band or two-margin
conditions may be used instead. Define
\[
 b(w)=D\{x:\mathsf{Stable}(w,x)\text{ fails}\}.
\]
Draw an independent unlabeled sample $U=(U_1,\ldots,U_m)$ and a
labeled sample $S=((X_i,f(X_i)))_{i=1}^n$. Condition
the raw Gaussian on stability at every input in $U$ and $S$, then
apply the usual likelihood $f(X_i)z_w(X_i)>1/2$.
Thus the support restriction depends on input locations but not on
their labels. Optional global clean-margin conditioning, as defined
after Theorem~\ref{thm:standalone-clipped},
may also be retained. Write $Q_{U,S}$ for the resulting posterior.

\begin{lemma}[Uniform coverage from an unlabeled stability pool]
\label{lem:unlabeled-stability-pacbayes}
With probability at least $1-\delta_U$ over $U$, every probability
law $Q$ supported on stability at all $U_j$ satisfies
\[
 \mathbb E_Q b(w)
 \le \frac{2}{m}\left[D_{\mathrm{KL}}(Q\Vert\mu)
                       +\log(1/\delta_U)\right].
\]
The assertion includes laws selected after observing the labeled
sample and its labels.
\end{lemma}

\begin{proof}
For a fixed $w$, the probability of stability at all $U_j$ is
$(1-b(w))^m$. Consequently
\[
 \mathbb E_U\int
 e^{m b(w)/2}\mathbf1\{\mathsf{Stable}(w,U_j)\text{ for every }j\}
 \,d\mu(w)
 \le\int e^{m b(w)/2}(1-b(w))^m\,d\mu(w)\le1.
\]
Markov's inequality bounds the displayed integral by $1/\delta_U$
on an event of probability at least $1-\delta_U$. On this event apply
the entropy variational inequality to the measure with density
$e^{m b/2}$ on the empirical stability set and zero outside it.
Since $Q$ is supported on that set, it gives
\[
 \frac m2\mathbb E_Qb
 \le D_{\mathrm{KL}}(Q\Vert\mu)+\log(1/\delta_U).
\]
Equivalently one may first truncate the logarithmic density below
and then pass to the limit. This proves the simultaneous assertion.
\end{proof}

\begin{theorem}[Width-independent labeled complexity with an unlabeled pool]
\label{thm:unlabeled-stability-pool}
Under the Boolean-target assumptions of the main theorem, let $D_*$
be its polynomial target-cylinder cost, so that the raw target cylinder
has mass at least $e^{-aD_*}$. For every fixed $N\ge C$ and
$0<t\le1$, there is a labeled sample bound polynomial in the same
parameters as the main theorem, independent of $N$ and $t$, such that
the following holds. If
\[
 m\ge\frac8\varepsilon
       \left[aD_*+\log(2/\delta)\right],
\]
then, with probability at least $1-\delta$ over $U,S$, the posterior
$Q_{U,S}$ has Gibbs test-margin error at most $\varepsilon$.
The usual posterior-confidence consequences follow.
\end{theorem}

\begin{proof}
Every member of the target cylinder is stable on the entire cube
and has the correct signed margin greater than one. Thus it belongs
to every conditioning event here, for every realization of $U,S$.
The posterior is a raw Gaussian conditioned on an event of probability
at least $e^{-aD_*}$, and therefore
\[
 D_{\mathrm{KL}}(Q_{U,S}\Vert\mu)\le aD_*.
\]
Lemma~\ref{lem:unlabeled-stability-pacbayes}, with
$\delta_U=\delta/2$, gives
$\mathbb E_{Q_{U,S}}b(w)\le\varepsilon/4$ under the displayed
unlabeled sample bound.

The same entropy bound gives the main theorem's readout truncation
and width-independent compressor. Choose the discarded posterior
mass at most $\beta=\varepsilon/4$ and the coupling failure probability
$\alpha=\varepsilon/12$. All posterior draws are stable at the labeled
input locations, so the finite candidate law has empirical ramp loss
at most $\alpha$. The same uniform finite-class event over $S$ bounds
its population ramp loss by
$2\alpha+8[\log|\mathcal H|+\log(2/\delta)]/n$.

At a fresh test input, loss transfer can fail only on a compression
failure or when the original network is unstable at that input.
Restoring the truncated mass and weighting the conditional instability
term by its retained mass gives
\[
 R(Q_{U,S})\le
 \beta+3\alpha+
 \frac{8[\log|\mathcal H|+\log(2/\delta)]}{n}
 +\mathbb E_{Q_{U,S}}b(w).
\]
Choose $n$ so that the finite-class term is at most $\varepsilon/4$.
The four terms then sum to $\varepsilon$. The compression class has
the same polynomial logarithmic size as before, independently of
$N,t,U,S$. The two sample events have total failure probability at
most $\delta$.
\end{proof}

This is an input-adapted prior with the usual label likelihood. In
particular, observing each labeled input also imposes an additional
label-free stability constraint; this is a specified change to the
learning rule, not a theorem about the original fixed prior. It
requires stability only at finitely many specified inputs, so it does
not infer region-wide stability from pointwise certificates. It also
does not provide a width-independent bound on the total data used:
at $t=N^{(1-\gamma)/2}$ its unlabeled cost contains $a=N^\gamma$.
For each fixed learner the probability statement above is complete.
To cover a prescribed countable collection of widths and scales
simultaneously, use the one common finite-class event for the labeled
sample with failure probability $\delta/2$, and allocate unlabeled
failure probabilities $\delta_j$ with $\sum_j\delta_j\le\delta/2$.
Use unlabeled sizes
$m_j\ge8[a_jD_*+\log(1/\delta_j)]/\varepsilon$; the labeled sample
size is unchanged, and the unlabeled pools may be nested prefixes
of one iid sequence. No
simultaneous statement over a continuum of changing Gaussian priors
is asserted without further argument.

An equivalent semi-supervised protocol chooses the prior before any
labels are revealed. Draw one iid pool of $M\ge\max\{m,n\}$ input
locations, impose stability on the whole pool, and then reveal the
labels of its first $n$ points. The coverage argument uses all $M$
inputs, and the finite-class generalization event uses the first $n$.
These two events need not be independent: each has its stated marginal
probability, and their failures are combined by a union bound. The same
guarantee therefore holds, with $n$ labels and $M$ input locations.
